% Matriz de consistencia (Guía N°2 EPG-UNI, p. 27) y matriz de trazabilidad.
% Compilar: latexmk -xelatex matriz_consistencia.tex
% Fuente del contenido: tesis, Capítulo I (§1.2–1.4) y Capítulo III (§3.3, Tabla de contrastación).
% Las celdas de la Tabla A marcadas "AUTORES" deben revisarse y redactarse con palabras propias
% antes de la entrega (docs/matriz_consistencia.md, sección 3).
\documentclass[10pt]{article}
\usepackage[a4paper,landscape,margin=1.5cm]{geometry}
\usepackage{fontspec}
\setmainfont{Arial}
\usepackage[spanish,es-noshorthands]{babel}
\usepackage{array,longtable,multirow}
\usepackage[table]{xcolor}
\usepackage{url}
\def\UrlBreaks{\do/\do_\do.}

\newcolumntype{L}[1]{>{\raggedright\arraybackslash}p{#1}}
\newcommand{\ingles}[1]{\textit{#1}}
\newcommand{\sub}[1]{\textbf{#1}\par}
\setlength{\parindent}{0pt}
\renewcommand{\arraystretch}{1.15}
\pagestyle{plain}

\newcommand{\titulo}{Sistema de análisis satelital de inundaciones basado en modelos de
segmentación profunda para el mapeo de áreas inundadas en el Perú}
\newcommand{\autores}{Gerson Andre Garrido Mansilla, Renzo Elvis Cerron Tome}

\begin{document}

{\centering
{\large\textbf{MATRIZ DE CONSISTENCIA}}\par\smallskip
Maestría en Ciencias con mención en Inteligencia Artificial, Universidad Nacional de Ingeniería\par}
\bigskip

% ============================================================================
\textbf{Tabla A.} Matriz de consistencia
\smallskip

{\scriptsize\setlength{\tabcolsep}{3.5pt}
\begin{longtable}{|L{2.9cm}|L{4.0cm}|L{4.0cm}|L{4.6cm}|L{3.9cm}|L{5.0cm}|}
\hline
\multicolumn{6}{|L{25.75cm}|}{\textbf{Título:} \titulo} \\ \hline
\textbf{Problema de estudio} & \textbf{Problemas} & \textbf{Objetivos} & \textbf{Hipótesis} &
\textbf{Variables y dimensiones} & \textbf{Diseño metodológico} \\ \hline
\endhead
% AUTORES: todas las celdas de esta tabla.
En el Perú no se cuenta con un sistema que produzca de forma automática mapas confiables de la
extensión de las inundaciones con imágenes Sentinel-2.

Los modelos entrenados con datos de otros países pierden desempeño cuando se aplican en la costa
norte peruana, y no se sabe qué tan confiable es la incertidumbre que reportan.

Para entrenar un modelo propio se necesitan etiquetas locales, que no suelen estar disponibles
después de un evento.
&
\sub{Problema general}
¿De qué manera se puede desarrollar un sistema de análisis satelital de inundaciones para el Perú,
basado en modelos de segmentación profunda, que genere mapas de la extensión del agua y de su
incertidumbre a partir de imágenes Sentinel-2, aun cuando no se dispone de datos etiquetados del
Perú para entrenar los modelos?
&
\sub{Objetivo general}
Desarrollar un sistema de análisis satelital de inundaciones para el Perú, basado en modelos de
segmentación profunda, que genere mapas de la extensión del agua y de su incertidumbre a partir de
imágenes Sentinel-2.
&
\sub{Hipótesis general}
Los modelos fundacionales geoespaciales pierden menos desempeño que la U-Net al pasar de
Sen1Floods11 al Perú.

La adaptación de dominio sin datos etiquetados del Perú y la calibración de la incertidumbre
mejoran en parte la confiabilidad de los mapas.

El sistema que integra estos modelos cumple sus requisitos funcionales y procesa una escena en un
tiempo adecuado para la respuesta después de un evento.
&
\sub{Variable independiente (X)}
Modelo de segmentación, estrategia de entrenamiento o adaptación y configuración del sistema.

\textbf{Dimensiones:}
X1. Arquitectura.\newline
X2. Esquema de entrenamiento.\newline
X3. Adaptación de dominio sin datos etiquetados del Perú.\newline
X4. Tratamiento de la incertidumbre.\newline
X5. Configuración del sistema.
&
\textbf{Tipo:} aplicada (desarrollo tecnológico).\newline
\textbf{Enfoque:} cuantitativo.\newline
\textbf{Nivel:} explicativo.\newline
\textbf{Método:} hipotético-deductivo; desarrollo iterativo e incremental para el sistema.\newline
\textbf{Diseño:} experimental (experimentos computacionales controlados) para los objetivos
específicos 1 y 2; evaluación técnica frente a requisitos fijados de antemano para el objetivo
específico 3.
\\ \hline
&
\sub{Problemas específicos}
PE1. ¿Qué modelo de segmentación (TerraMind, Prithvi o U-Net) y qué esquema de entrenamiento
logran el mejor desempeño en Sen1Floods11 y en el Perú?

PE2. ¿En qué medida la adaptación de dominio sin datos etiquetados del Perú reduce la pérdida de
desempeño de los modelos en el Perú, y qué tan bien calibrada está su incertidumbre?

PE3. ¿En qué medida el sistema implementado cumple los requisitos funcionales y el tiempo de
procesamiento necesarios para usarlo después de una inundación?
&
\sub{Objetivos específicos}
OE1. Comparar el IoU de agua de TerraMind v1-B, Prithvi-EO-2.0 300M y U-Net en Sen1Floods11 y en
los eventos peruanos de 2023 y 2017, con tres esquemas de entrenamiento, para seleccionar el
modelo del sistema.

OE2. Determinar en qué medida la adaptación de dominio no supervisada y la calibración de la
incertidumbre mejoran la confiabilidad de los mapas en el Perú, e identificar las condiciones en
las que se concentran los errores.

OE3. Implementar el sistema como una aplicación web y evaluar si cumple los requisitos funcionales
y el tiempo de procesamiento por escena Sentinel-2.
&
\sub{Hipótesis específicas}
HE1. En Sen1Floods11 no se detecta una diferencia significativa de IoU de agua entre los modelos
fundacionales y la U-Net (IC 95~\%). En el Perú, la caída de IoU de los modelos fundacionales es
menor que la de la U-Net.

HE2. Los errores en el Perú se concentran en pocas condiciones: bordes de nubes, vegetación
inundada y terreno desértico o costero. La adaptación aumenta el IoU de agua respecto del modelo
sin adaptar (IC 95~\% mayor que cero), sin llegar al modelo entrenado en el Perú. Sin calibrar, los
modelos son sobreconfiados; el escalado de temperatura y los ensambles reducen el ECE, y la
incertidumbre permite detectar errores (AUROC mayor que 0.5).

HE3. El sistema cumple el 100~\% de los requisitos funcionales RF1 a RF7 y procesa una escena
Sentinel-2 completa en 15 minutos o menos, incluida la descarga.
&
\sub{Variable dependiente (Y)}
Calidad del mapeo de inundaciones.

\textbf{Dimensiones:}
Y1. Exactitud de la segmentación.\newline
Y2. Generalización al Perú.\newline
Y3. Calibración.\newline
Y4. Detección de errores.\newline
Y5. Desempeño del sistema.

\sub{Variables intervinientes}
Evento (2023 o 2017), calidad de las etiquetas y cobertura de nubes.
&
\textbf{Población:} imágenes Sentinel-2 L1C de eventos de inundación.\newline
\textbf{Muestra (no probabilística):} Sen1Floods11, particiones oficiales (252 de entrenamiento,
89 de validación y 90 de prueba); Perú 2023, 1\,291 recortes por bloques espaciales (753, 174 y
172, más 192 de separación sin uso); Perú 2017, 204 recortes solo de prueba (101 en la evaluación
principal); escenas completas para el sistema.\newline
\textbf{Técnicas e instrumentos de recolección de datos:} descarga desde Copernicus Data Space;
etiquetas generadas por reglas a partir de los mapas de UNOSAT (2023) y Copernicus EMS (2017),
sin etiquetado manual; programas del repositorio versionados con \ingles{git}; registros de
pruebas y de tiempos del sistema.\newline
\textbf{Técnicas e instrumentos de análisis y procesamiento:} IoU de agua, precisión,
\ingles{recall} y mIoU; ECE; AUROC; diferencias pareadas con IC 95~\% por remuestreo
\ingles{bootstrap} (2\,000 réplicas); media y desviación estándar de tres semillas; lista de
verificación de RF1 a RF7; tiempo por escena.
\\ \hline
\end{longtable}}

Autores: \autores.\\
Elaboración: propia, a partir de los Capítulos I y III de la tesis.

\newpage
% ============================================================================
\textbf{Tabla B.} Matriz de trazabilidad: de cada hipótesis específica a su indicador, instrumento,
evidencia y estado
\smallskip

{\scriptsize\setlength{\tabcolsep}{3.5pt}
\begin{longtable}{|L{0.8cm}|L{1.3cm}|L{2.2cm}|L{4.1cm}|L{3.2cm}|L{3.8cm}|L{4.0cm}|L{5.0cm}|}
\hline
\textbf{Afirm.} & \textbf{PE / OE / HE} & \textbf{Dimensión} & \textbf{Indicador} &
\textbf{Criterio de contrastación} & \textbf{Instrumento} & \textbf{Evidencia} & \textbf{Estado} \\
\hline
\endhead
1a & PE1 / OE1 / HE1 & Y1. Exactitud de la segmentación &
Diferencia pareada de IoU de agua entre la U-Net y cada modelo fundacional, prueba de Sen1Floods11 &
El IC 95~\% incluye el cero &
\path{scripts/evaluate_frozen.py}, \path{scripts/bootstrap_frozen.py} &
\path{outputs/phaseD/test/bootstrap_diffs.csv}; Cap.~III, §3.3 &
Coherente: no se detectó una diferencia. U-Net $-$ Prithvi: +0.013 [$-$0.022; 0.062]; U-Net $-$
TerraMind: +0.015 [$-$0.018; 0.051]. \\ \hline
1b & PE1 / OE1 / HE1 & Y2. Generalización al Perú &
Diferencia de caídas $D$ = caída de la U-Net $-$ caída del modelo fundacional (Sen1Floods11 $\to$
prueba Perú 2023) &
IC 95~\% de $D$ por encima de cero &
\path{scripts/bootstrap_drop_diff.py} (análisis \ingles{post hoc}) &
\path{outputs/phaseD/test/bootstrap_drop_diff.csv} &
Se confirma parcialmente. Prithvi: $D$ = 0.177 [0.102; 0.253]. TerraMind: $D$ = 0.064 [0.001;
0.120], con semillas solapadas con la U-Net. \\ \hline
2a & PE2 / OE2 / HE2 & Y1. Exactitud de la segmentación &
Razón de concentración por estrato (proporción de errores / proporción de píxeles), validación
peruana &
Razón $\geq$ 2 en los estratos nombrados (criterio fijado después de ver los resultados) &
\path{scripts/error_shares_dev.py} &
\path{outputs/phaseD/analysis/error_shares.csv} &
Respaldo exploratorio parcial: solo la vegetación inundada cumple el criterio (U-Net y Prithvi);
la mayor concentración está en el borde del agua. \\ \hline
2b & PE2 / OE2 / HE2 & Y2. Generalización al Perú &
Diferencia pareada de IoU de agua: adaptado $-$ origen y adaptado $-$ entrenado en el Perú &
IC 95~\% de la primera $>$ 0 y del modelo adaptado por debajo del entrenado en el Perú &
\path{scripts/evaluate_frozen.py}, \path{scripts/bootstrap_frozen.py} &
\path{outputs/phaseD/test/bootstrap_diffs.csv} &
Se confirma en 2023: U-Net +0.222 [0.132; 0.309]; TerraMind +0.082 [0.016; 0.156]. Se refuta frente
a la referencia de 2017: U-Net $-$0.063; TerraMind $-$0.103. Prithvi sin configuración
adaptada. \\ \hline
2c & PE2 / OE2 / HE2 & Y3. Calibración; Y4. Detección de errores &
Temperatura ajustada $T$; ECE (\%); AUROC de la incertidumbre &
$T > 1$; ECE menor con temperatura y con ensamble; AUROC $>$ 0.5 &
\path{scripts/phaseD_analysis.py}, \path{scripts/calibration_frozen.py} &
\path{outputs/phaseD/analysis/temperatures.json}; \path{outputs/phaseD/test/calibration.md} &
Se confirma en las configuraciones congeladas: AUROC de 0.598 a 0.783. Excepción fuera de ellas:
Prithvi entrenado en el Perú (0.439). \\ \hline
3 & PE3 / OE3 / HE3 & Y5. Desempeño del sistema &
Requisitos funcionales cumplidos (\%); tiempo por escena (min), 6 teselas $\times$ 3 repeticiones &
100~\% de RF1 a RF7; mediana $\leq$ 15 min en cada tesela &
\path{scripts/system_rf*_check.py}, \path{scripts/system_timing.py};
\path{docs/system_evaluation_protocol.md} &
\path{outputs/system_eval/} &
Pendiente de la medición formal. Verificaciones preliminares superadas: RF1, RF3 y RF5. \\ \hline
\end{longtable}}

\textit{Nota.} Los intervalos son IC 95~\% por remuestreo \ingles{bootstrap} pareado de recortes
(2\,000 réplicas) sobre el IoU medio de tres semillas. Las cifras de 2017 corresponden a los 101
recortes que no se superponen con el entrenamiento de 2023. Las etiquetas de 2017 conservan solo
los píxeles en que EMS y el índice MNDWI concuerdan, por lo que miden sobre todo la exhaustividad.
El detalle de cada veredicto está en el Capítulo III, §3.3.

Autores: \autores.\\
Elaboración: propia.

\end{document}
